\chapter{有机反应不必一条条死背}

\begin{kaichang}
请你到厨房里找三样东西：一个苹果、一瓶医用消毒酒精、一袋用塑料包装袋装着的面包。

现在把苹果切开，放上半个小时。你会看到雪白的果肉慢慢变成难看的茶色。与此同时，那瓶酒精静静待着，那袋面包也静静待着——看起来什么都没发生。

先猜三个问题：苹果变褐、酒精能燃烧、塑料袋是用什么办法“造”出来的——这三件毫不相干的事，背后有没有一套共同的语言？如果有一本写满几千个有机反应的“反应词典”摆在你面前，你觉得必须背下多少条，才能算学会有机化学？

把答案先写在纸上。这一章结束时我们回来对答案。我提前剧透一句：需要“背”的反应，可能比你想象的少得多。
\end{kaichang}

\section{几千个反应，一双手数得过来的套路}

翻开任何一本大学有机化学教材，你都会被吓到：一页接一页的反应式，动辄上千个。19 世纪的化学家们确实是这么过来的——每发现一个新反应，就给它起个名字，记在本子上，像集邮一样。

先算一笔账，看看“集邮式学习”有多绝望。今天人类已知的含碳化合物超过三千万种，而且每天还在增加。就算每个化合物只对应十个常见反应，那也是三亿条反应式。假如你记忆力超群，一天能背五十条，风雨无阻地背，需要一万六千年。显然，没有任何一个化学家是这么学的。他们用的是另一套办法：\textbf{把成千上万的反应归进少数几个家族}。

其实你自己早就是有机反应的目击证人了。香蕉放几天由青转黄，是果实内部的淀粉被拆解、风味分子被改写；牛奶放酸了，是乳糖被细菌改造成了乳酸；烤肉表面那层焦香，是糖和氨基酸在高温下牵手上百次的结果；连你的头发能被烫发剂改变形状，靠的也是打断又重接蛋白质里的二硫键。这些变化的名字千奇百怪，可如果把镜头推近到原子层面，它们反复使用的只是同一套动作。

归类的线索，你在前面的章节也已经攒齐了。第 36 章告诉你碳骨架长什么样，第 37 章告诉你官能团是分子身上的“反应接口”——羟基、羰基、羧基、碳碳双键，接口不同，脾气不同。这一章要解决的是最后一个问题：这些接口上发生的反应，有没有通用的剧本？

答案是：有。而且剧本的核心角色只有三种，全部来自你熟悉的老朋友——电子。

\section{先找到“电子多”和“电子缺”的位置}

第 11 章讲电负性时说过：两个原子成键，如果电负性不同，共享的电子对就会偏向电负性大的那一方。这不是口号，它是理解几乎所有有机反应的第一块基石。

看一个具体例子。丙酮——洗甲水里的主要成分——分子中央有一个碳氧双键（\chem{C{=}O}，叫羰基）。氧的电负性比碳大得多，于是双键里的电子云被拉向氧那一边。结果是什么？氧原子带上少量负电荷，成了\textbf{电子丰富}的位置；而羰基碳则\textbf{电子贫乏}，像个没吃饱饭的家伙。整根键就像一个两头温度不同的房间：一头挤，一头空。

再想想这个图景意味着什么。如果这时飘来另一个电子丰富的分子碎片，它会被哪儿吸引？同性相斥，它绝不会去碰电子已经很多的氧；它会直奔那个缺电子的羰基碳。反过来，一个缺电子的正离子来了，它会去找氧。于是，你甚至不用知道反应叫什么名字，只要看清分子上哪里电子多、哪里电子缺，就能大致猜出反应会从哪里开打。

化学家给这种“电子分布不均匀”准备了一个记号：在电子偏少的一端标 $\delta+$（读作“德尔塔正”），在电子偏多的一端标 $\delta-$。这不是完整的电荷，只是“略微带一点电性”的意思，但已经足够预测行为。再练一个：溴乙烷里的碳溴键（\chem{C{-}Br}），溴的电负性更大，所以溴是 $\delta-$、碳是 $\delta+$——这个缺电子的碳，就是溴乙烷最容易被“攻打”的城门，稍后讲取代反应时你就明白为什么。

化学家给这两类角色起了名字，名字唬人，意思却朴素：

\begin{itemize}[itemsep=2pt]
\item \textbf{亲核体}：电子丰富、愿意拿出一对电子去成键的家伙。带负电的离子（如氢氧根 \chem{OH^-}）常常是亲核体，带孤对电子的中性分子（如水、氨）也是。“亲核”就是“喜欢靠近带正电的原子核”。
\item \textbf{亲电体}：电子贫乏、想从别人那里接过一对电子的家伙。刚说的羰基碳、氢离子 \chem{H^+}，都是亲电体。
\item \textbf{自由基}：第 33 章讲氧化还原时提过，电子可以一个一个地转移。当一根共价键被“劈开”、两个原子各分走一个电子时，就产生了带着单电子的碎片——自由基。它不守“电子对”的规矩，像击鼓传花一样把单电子挨个传下去。燃烧、塑料老化、防晒霜要防的皮肤损伤，背后都有它，第 39 章还会专门讲它。
\end{itemize}

现在你可以把大多数有机反应概括成一句话了：\textbf{电子丰富的位置去找电子贫乏的位置，电子对搬一次家，旧键断开，新键生成。}分子没有意愿，不会“想”反应；只是电子对从高处（能量高、受约束少的位置）流向低处时，整个体系的总能量降低了，变化就可能发生——这正是第 27、28 章能量账本在有机分子身上的又一次上演。

\section{弯箭头：画的是电子，不是原子}

剧本有了，还得有一套记谱法。有机化学家的记谱法叫\textbf{弯箭头}，规矩只有两条，却常被误解：

\begin{itemize}[itemsep=2pt]
\item 一条普通的弯箭头（箭头是完整的尖）代表\textbf{一对}电子的搬家：箭尾从电子对现在的位置出发，箭头指向电子对要去的位置。
\item 一条“鱼钩”式的半箭头（只有半个尖）代表\textbf{一个}电子的移动，专门用于自由基。
\end{itemize}

拿一个最简单的反应来试。把乙烯（\chem{CH_2{=}CH_2}，水果催熟时释放的气体）通进溴化氢（\chem{HBr}）里，会生成溴乙烷：

\[\chem{CH_2{=}CH_2} + \chem{HBr} \rightarrow \chem{CH_3CH_2Br}\]

用弯箭头的语言讲，这件事分两步记谱：第一步，乙烯双键里那对不太安分的 $\pi$ 电子，画一条弯箭头，指向 \chem{HBr} 中带部分正电的氢；同时再画一条弯箭头，让 \chem{H{-}Br} 键里的电子对整个搬到溴身上。第二步，带着负电和孤对电子的溴离子（现在它是亲核体了）画一条弯箭头，指向刚形成的、缺电子的碳，补上最后一根键。

请注意这里最容易滑倒的地方：\textbf{弯箭头记录的是电子的账本，不是原子的飞行路线图}。真实的分子在溶液里翻滚、碰撞、振动，场面混乱得多，没有任何一个原子沿着你画的那条弧线飞行。弯箭头只回答一个问题：哪对电子从哪儿搬到了哪儿。可一旦这份电子账本记清楚了，原子最终的新位置、新键断旧键的结局，就全都能推出来——这就是它的威力。

\section{七大家族：取代、加成、消去、氧化、还原、缩合、水解}

把“电子搬家”的剧本按搬家方式分分类，几千个有机反应就被收进了七个家族。每个家族我都配一个你见过的例子。

\subsection{取代：一个换一个}

分子里某个位置上的原子或原子团，被另一个原子团顶掉。亲核体拿着电子对来“顶岗”，原来的住户带着电子对离开。医院缝合伤口用的某些可吸收缝线、你稍后会在第 40 章细看的酯化反应，本质上都是取代。漂白剂在光照下和有机物的反应，走的也是取代的路——所以彩色衣服沾上漂白水会花掉。

\subsection{加成：双键打开，两边各加一个}

碳碳双键里那对 $\pi$ 电子是现成的亲核弹药，遇到亲电体就打开双键，两头各接上一个新原子。刚才的乙烯加溴化氢就是。食品工业里把液态植物油变成半固态的人造奶油，用的“氢化”工艺就是往双键上加氢——这也是配料表上“氢化植物油”四个字的来历。

\subsection{消去：加成的倒放}

从一个分子上“剪掉”一小段（常是水或卤化氢），让相邻两个碳之间长出双键。它是加成的逆过程，就像同一个剧本倒着演。工业上大量制造乙烯，靠的正是从更大的烃分子上消去氢气。

\subsection{氧化与还原：老朋友换了件有机的外衣}

第 33 章已经说透：氧化还原不等于“得氧失氧”，本质是电子的转移和氧化数的升降。在有机化学里有个好用的速记：分子\textbf{加氧或失氢}，多半是氧化了；\textbf{加氢或失氧}，多半是被还原了。开场那个苹果变褐，就是果肉里的酚类物质被空气中的氧气氧化；食用油放久了产生哈喇味，是油脂被氧慢慢氧化的结果；而刚才说的植物油加氢，从另一头看，也正是双键被还原了。

\subsection{缩合与水解：手拉手与剪一刀}

两个分子“手拉手”连成一个大分子，同时挤出一小分子水（或其他小分子），叫\textbf{缩合}；反过来，用水把一个大分子“剪”成两段，叫\textbf{水解}。这一对是生命中的头号反应：氨基酸缩合成蛋白质、葡萄糖缩合成淀粉，而你吃饭后的消化，几乎全靠水解把淀粉、蛋白质剪回小分子——第六篇讲生物大分子时，你会和它们反复见面。走出生命世界，这对反应同样忙碌：尼龙、涤纶这些合成纤维，都是小分子一次次缩合拉成的长链（第 43 章细讲）；做肥皂的皂化反应，是油脂的水解，第 40 章会带着香气和酸味细讲它。

再强调一次：这七大家族不是法律条文，不同教材的划分略有出入。它们是七副眼镜，帮你在一团乱麻的反应现象里看出秩序。以后每遇到一个新反应，先问一句：这是电子对在取代、加成、消去，还是在氧化还原、缩合水解？十有八九，你能立刻给它找到家谱。

\begin{qiaoliang}
同一个反应物，有时能生成两种甚至更多种产物。往哪边走，由两本账共同决定，而这两本账你在第一册已经见过：哪条路\textbf{门槛}低（活化能，第 29 章），决定哪条走得快——这是\textbf{动力学}；哪个\textbf{终点}低（自由能 $\Delta G$，第 28 章），决定最后停在哪——这是\textbf{热力学}。

经典例子：1,3-丁二烯和 \chem{HBr} 加成，可以生成双键在中间的产物（能量更低、更稳定），也可以生成双键在端头的产物（生成的门槛更低、速度更快）。实验发现：低温、短时间反应，端头产物占上风——大家只来得及翻过矮墙，这叫\textbf{动力学控制}；升高温度、延长反应时间，产物之间能来回转化，最终大部分都汇聚到能量最低的那个谷底，这叫\textbf{热力学控制}。

\begin{center}
\begin{tikzpicture}[scale=0.9]
\draw[->] (0,0) -- (8.4,0) node[right]{反应进程};
\draw[->] (0,0) -- (0,4.2) node[above]{能量};
\draw (1,1.4) .. controls (1.7,2.1) and (2.1,2.6) .. (2.6,2.8)
      .. controls (3.1,3.0) and (3.4,2.2) .. (4,0.9);
\draw (1,1.4) .. controls (1.8,2.9) and (2.3,3.5) .. (2.9,3.6)
      .. controls (3.5,3.7) and (3.9,1.6) .. (4.6,0.3);
\node at (0.8,1.1) [below left]{反应物};
\node at (4.2,0.9) [above right]{产物甲（快，门槛矮）};
\node at (4.9,0.2) [below right]{产物乙（稳，谷底深）};
\end{tikzpicture}
\end{center}

（这是人为的能量示意图，曲线形状不代表精确测量。）一句话记住：\textbf{动力学管“多快到”，热力学管“停在哪”}。把自发当成快、把放热当成必然，是把两本账混成了一本——那是第 28、29 章就拆穿过的误会。
\end{qiaoliang}

\section{机理不是口诀，是被证据钉住的模型}

到这里，你可能会产生一种不安：电子对这样搬、那样搬，讲得头头是道——可谁也没有看见过电子搬家。这些弯箭头，会不会只是化学家编出来的顺口溜？

这个质疑非常正当，而化学家自己也早就想到了。他们给“反应一步步怎样发生”的描述起了个专门的名字：\textbf{反应机理}。机理是模型，不是实况录像；但一个合格的机理必须被实验反复检验，能解释旧事实，还能预言新现象。下面这个故事，就是机理怎样被钉上证据钉子的经典案例。

\begin{zhengju}
乙酸乙酯水解，生成乙酸和乙醇：

\[\chem{CH_3COOCH_2CH_3} + \chem{H_2O} \rightarrow \chem{CH_3COOH} + \chem{CH_3CH_2OH}\]

方程式两边看起来天经地义，但机理上有一个死结：酯分子中间那根“碳—氧”键有两种断法。水解时，是水把酰基那一侧的键剪断（乙酸拿走水里的氧），还是把醇那一侧的键剪断（乙醇拿走水里的氧）？两种断法写出的总方程式一模一样，光看方程式根本分不出来。换句话说：产物相同，机理却可以完全不同——这正是“方程式不是机理”的活教材。

突破口来自氧的一种特殊版本。天然氧绝大多数是氧-16，但有极少量的氧-18——原子核里多了两个中子，化学脾气完全一样，却能被质谱仪（第 35 章）区分开。20 世纪 30 年代末到 40 年代，化学家（如波拉尼与萨博等人的工作）用\textbf{氧-18 标记的水}去水解普通的酯：如果乙酸那侧断键，标记氧应该跑进乙酸；如果乙醇那侧断键，标记氧就该出现在乙醇里。

实验结果干净利落：氧-18 进了乙酸，乙醇里几乎找不到。于是机理被钉死了：断的是酰基和氧之间的键，水里的氧最终归属于酸。今天教材上酯水解的弯箭头画法，就是这样画出来的——每一条箭头背后，都站着同位素标记、反应速率测定、中间体捕捉这一整套证据。顺便说，这个方法后来威力更大：光合作用放出的氧气来自水还是二氧化碳这个世纪之问，也是用同样的氧-18 标记思路解决的（第六篇会讲）。机理是模型，但它是\textbf{戴着证据镣铐跳舞的模型}——新证据来了，它就得改。
\end{zhengju}

\begin{shiyan}
\textbf{苹果变褐大侦探}（在家就能做，安全）

\textbf{材料}：一个苹果、两个干净小碟、半个柠檬（或一片维生素 C 泡腾片化在水里）、小刀。

\textbf{步骤}：请大人把苹果切成四片。两片放在碟子 A 里，表面什么都不涂；两片放在碟子 B 里，表面均匀涂上柠檬汁（或维生素 C 溶液）。都摆在桌上，不要盖。

\textbf{观察点}：30 分钟、1 小时、2 小时后各看一次，比较两个碟子里苹果切面的颜色深浅。可以再闻闻气味有没有差别。

\textbf{安全提示}：切苹果请大人代劳；实验后的苹果暴露在空气里太久，不建议再吃——不是因为它变“毒”了，而是敞放的切面容易沾上微生物。

你会看到：涂了柠檬汁的苹果明显变褐得慢。为什么？变褐是果肉中的酚类物质被空气中的氧气\textbf{氧化}（氧化家族）。柠檬汁里有双重防线：维生素 C 是个积极的还原剂，抢在酚类前面“替它们被氧化”；酸性的环境又让催化这一步的酶放慢手脚。一个厨房里的小对比，同时用上了氧化还原、催化剂和第 29 章的速率知识。
\end{shiyan}

\section{这张地图的边界在哪里}

七大家族加弯箭头，是一张非常好用的地图。但地图不是领土，它有三条明确的边界，值得你记在页边空白处。

第一，\textbf{同一个分子可能有好几个“接口”}，往哪儿反应取决于条件。溶剂、温度、浓度、酸碱环境一变，占上风的产物可能就换了——热力学控制和动力学控制的拔河，就是最常见的例子。所以“背产物”不如“看条件”。

第二，\textbf{有些反应不走亲核—亲电的剧本}。有一类反应里，几根键同时断开、同时形成，电子像围成一圈接力，没有明确的亲核体和亲电体之分。你不需要现在掌握它，只要知道：七大家族覆盖了绝大多数常见反应，但不是全部。

第三，\textbf{机理本身会被修订}。氧-18 的故事告诉你机理靠证据立足；反过来，新证据也可能推翻旧机理。有机化学教科书每隔若干年就会修订一些机理的细节。所以正确的心态是：把机理当作“目前证据下最好的解释”，而不是刻进石头的教义。这也是“不必死背”的深层原因——要背的是判据，不是结论。

\begin{wuqu}
“弯箭头画的是原子从一个位置飞到另一个位置的路线。”——错。弯箭头记账的对象是\textbf{电子对}，不是原子。真实碰撞中的原子并不沿那条弧线运动；甚至有这样的事：亲核取代里弯箭头从亲核体指向碳，而被顶掉的基团是从\textbf{相反一侧}离开的——实验上表现为分子的空间构型整个翻转，像一把伞被大风从里吹到外（这叫瓦尔登翻转，1893 年就被观察到了）。如果箭头真是原子的飞行路线，这个“背面进出”的结果就无法解释。

另一个相关的误解是：“自由基是罕见的怪物。”恰恰相反。氧气分子自己就带着两个配对不上的单电子；你点火做饭的燃烧、塑料在太阳下发脆、铁生锈，全是自由基在击鼓传花。它不罕见，只是不守电子对的规矩。
\end{wuqu}

\section{回到厨房里的三样东西}

现在兑现开场的承诺。

\textbf{切开的苹果}：变褐是酚类物质的氧化反应——电子从酚流向氧，氧化数此升彼降。它属于七大家族里的氧化，而且被酶催化着（第 30 章说过，催化剂只降低门槛、不改总账）。柠檬汁能拖慢它，靠的是还原剂抢先被氧化。

\textbf{那瓶酒精}：它静静待着时什么也没发生，因为乙醇和氧气虽然“账面上”完全该反应（燃烧放出大量能量），反应的门槛却很高——这正是“自发不等于快速”的又一个实例，第 28、29 章的核心在这里又冒了头。划一根火柴提供越过活化能的能量，它就走自由基的路烧起来——那是第 39 章的主题。

\textbf{塑料包装袋}：它多半来自乙烯分子一个接一个的加成——双键打开，首尾相连，成千上万次重复之后，气体变成了薄膜。这是加成家族的“无限连播版”，第 43 章讲塑料、橡胶和纤维时会专门展开。

你看，三样东西、三种变化，没有一条是靠“背”出来的。极性告诉你电子在哪儿，弯箭头告诉你电子去哪儿，七大家族告诉你这是哪一类搬家——剩下的全是推理。这就是“不必死背”的真正含义：不是偷懒，而是把记忆力换成判断力。

\section{从“认套路”到“走路线”}

学会认套路之后，化学家做的事更进一层：\textbf{用套路组合出路线}。想造一种新药、一种新的可降解塑料，就从目标分子倒推：最后一步用缩合？倒数第二步用取代？像下棋一样一步步倒推到货架上买得到的原料——这叫逆合成思维，第 44 章会专门讲。而在那之前，我们先回到最古老、规模最大的那一类有机反应：把地下挖出来的黑色液体变成汽油、柴油和塑料原料。燃烧和裂化里主角不是电子对，而是击鼓传花的自由基；一锅石油怎样按沸点分开、又怎样被“剪碎”重新组合，里面还藏着一氧化碳为什么这么毒——第 39 章，我们从地下那锅黑汤讲起。

\begin{tiaozhan}
同样是“取代”，凭什么有的快、有的慢、有的还翻转构型？把速率当作侦探工具，能分出两种机理。以卤代烃被氢氧根取代为例：如果实验发现速率只取决于卤代烃的浓度、跟 \chem{OH^-} 的浓度无关，说明决定快慢的那一步只有卤代烃一个分子参与——它先自己断掉碳卤键、生成一个缺电子的碳正离子，亲核体随后才来补位（这叫单分子取代，SN1）。如果速率与两者的浓度都有关，说明关键一步是亲核体从背面“撞”上去、旧键新键同时交接（双分子取代，SN2）。这套用\textbf{速率方程反推机理}的方法，是 20 世纪 30 年代英戈尔德等人奠定的；它也是第 29 章“速率方程来自实验、不能从方程式猜”那条规则最漂亮的战果。跳过本框不影响主线。
\end{tiaozhan}

\section*{练一练}

\begin{sikaoti}
\item 丙酮的羰基（\chem{C{=}O}）和乙醇的羟基（\chem{O{-}H}）都含氧。分别标出这两个分子里电子丰富和电子贫乏的原子，并预测：一个带负电的亲核体靠近丙酮时，最可能进攻哪个原子？说明理由。
\item 给下列生活现象归家族（取代、加成、消去、氧化、还原、缩合、水解），并各用一句“电子搬家”的话解释：①敞口的食用油放久产生哈喇味；②用乙烯让青涩的香蕉变黄变甜（提示：乙烯是信号，催熟过程里淀粉被分解成糖）；③吃完米饭，淀粉最终变成葡萄糖被小肠吸收。
\item 用弯箭头的语言，把“乙烯 + \chem{HBr} $\rightarrow$ 溴乙烷”分成两步讲给同桌听：每一步哪对电子从哪里搬到哪里？再回答：这两条弯箭头代表溴原子和氢原子的飞行路线吗？
\item 某反应在低温下主要生成产物 A，升温并延长反应时间后主要生成更稳定的产物 B。画出类似本章能量示意图的草图，标出 A 和 B 的位置，并分别用“动力学控制”“热力学控制”解释两个结果。
\item 氧-18 标记实验：如果用普通水去水解一个“醇那一侧的氧被标记成氧-18”的酯，按照本章证实的机理，反应结束后标记的氧-18 会出现在乙酸里还是乙醇里？说说你的推理。
\item 找一篇“补充维生素 C 能永葆青春”之类的广告文案，用本章和第一册的氧化还原知识指出它哪些说法过度简化了“抗氧化剂”的作用（提示：剂量、证据等级、个体差异，以及人体内不是一个碟子里的苹果）。
\end{sikaoti}
